\documentclass[12pt,letterpaper]{article}

\usepackage[utf8]{inputenc}
\usepackage[T1]{fontenc}
\usepackage[spanish,es-tabla]{babel}
\usepackage{geometry}
\usepackage{xcolor}
\usepackage{hyperref}
\usepackage{longtable}
\usepackage{array}
\usepackage{booktabs}
\usepackage{enumitem}
\usepackage{listings}
\usepackage{fancyhdr}
\usepackage{titlesec}

\geometry{margin=2.4cm}
\hypersetup{
  colorlinks=true,
  linkcolor=blue!55!black,
  urlcolor=blue!55!black,
  pdftitle={Reporte tecnico del demo web},
  pdfauthor={Codex}
}

\definecolor{codebg}{RGB}{245,247,250}
\definecolor{codeframe}{RGB}{210,218,226}
\definecolor{heading}{RGB}{20,40,60}

\lstdefinestyle{reportcode}{
  backgroundcolor=\color{codebg},
  frame=single,
  rulecolor=\color{codeframe},
  basicstyle=\ttfamily\small,
  breaklines=true,
  columns=fullflexible,
  keepspaces=true,
  showstringspaces=false
}
\lstset{style=reportcode}

\titleformat{\section}{\Large\bfseries\color{heading}}{\thesection.}{0.6em}{}
\titleformat{\subsection}{\large\bfseries\color{heading}}{\thesubsection.}{0.6em}{}

\pagestyle{fancy}
\fancyhf{}
\lhead{Demo chatbot con agente orquestador}
\rhead{\thepage}
\renewcommand{\headrulewidth}{0.4pt}

\setlist[itemize]{leftmargin=1.2cm}
\setlist[enumerate]{leftmargin=1.2cm}
\renewcommand{\arraystretch}{1.2}

\begin{document}

\begin{titlepage}
  \centering
  \vspace*{2.0cm}
  {\Huge\bfseries Reporte tecnico\par}
  \vspace{0.4cm}
  {\Large Demo web de chatbot con agente orquestador\par}
  \vspace{0.6cm}
  {\large Detector de intenciones, refuerzo OpenIA offline y visualizacion de flujo de tokens\par}
  \vfill
  \begin{tabular}{rl}
    Proyecto: & EjemploApp \\
    Stack: & React + Vite + TypeScript / FastAPI + Pydantic \\
    Modalidad: & Demo local sin APIs externas \\
    Documento: & Reporte de entrega \\
  \end{tabular}
  \vfill
  {\large \today\par}
\end{titlepage}

\tableofcontents
\newpage

\section{Resumen ejecutivo}

Se desarrollo un demo web funcional de una aplicacion tipo chatbot con agente orquestador, detector de intenciones, refuerzo OpenIA/OpenAI-compatible en modo offline, reglas de ciberseguridad y visualizacion interna del flujo de tokens.

La aplicacion permite que un usuario escriba un mensaje, lo procese en backend por etapas y visualice en frontend el flujo completo mediante nodos secuenciales, inspector tecnico JSON, metricas de tokens, intencion detectada, agente seleccionado, herramienta ejecutada y reglas de seguridad aplicadas.

El demo corre localmente sin APIs externas, sin claves reales y sin ejecutar codigo arbitrario.

\section{Alcance implementado}

\begin{itemize}
  \item Frontend en React + Vite + TypeScript.
  \item Visualizacion de flujo con React Flow.
  \item Backend en Python + FastAPI.
  \item Validacion de requests y respuestas con Pydantic.
  \item Reglas de ciberseguridad en YAML.
  \item Pipeline trazable por etapas.
  \item Detector de intenciones por reglas.
  \item Modulo de refuerzo de intencion \texttt{OpenIA/OpenAI-compatible} en modo offline.
  \item Orquestador de agentes.
  \item Allowlist de herramientas locales.
  \item Ejemplos educativos de ataques de IA.
  \item Documentacion tecnica con diagramas y flujos.
\end{itemize}

\section{Estructura del proyecto}

\begin{lstlisting}[language={},caption={Estructura general del repositorio}]
EjemploApp/
  backend/
    app/
      agents/
      services/
      tools/
      main.py
      schemas.py
    security_rules.yaml
    requirements.txt
  frontend/
    src/
      components/
      services/
      styles/
      types/
      App.tsx
    package.json
    vite.config.ts
  docs/
    FUNCIONAMIENTO_Y_FLUJOS.md
    REPORTE_TECNICO.md
    REPORTE_TECNICO.tex
  README.md
\end{lstlisting}

\section{Arquitectura general}

La arquitectura se divide en frontend, API backend, servicios de negocio, agentes y herramientas locales.

\begin{center}
\begin{tabular}{>{\bfseries}p{4cm}p{10cm}}
\toprule
Componente & Responsabilidad \\
\midrule
Frontend React & Captura mensajes, renderiza chat, muestra metricas, visualiza nodos y permite inspeccionar etapas. \\
FastAPI & Expone endpoints HTTP y valida contratos con Pydantic. \\
Pipeline & Ejecuta las etapas secuenciales y genera trazas estructuradas. \\
Detector de intencion & Clasifica la solicitud inicial mediante reglas simples. \\
Refuerzo OpenIA & Ajusta la intencion en modo offline con senales semanticas locales. \\
Policy guard & Bloquea prompt injection, secretos y exceso de presupuesto. \\
Orquestador & Selecciona agente y herramienta segun la intencion final. \\
Agentes & Generan respuestas deterministas locales. \\
Herramientas & Ejecutan funciones permitidas por allowlist. \\
\bottomrule
\end{tabular}
\end{center}

\section{Backend}

El backend expone una API FastAPI con tres responsabilidades principales:

\begin{enumerate}
  \item Recibir mensajes del usuario.
  \item Ejecutar el pipeline interno.
  \item Devolver una respuesta estructurada con trazas, seguridad, tokens e intencion.
\end{enumerate}

\subsection{Endpoints}

\begin{longtable}{p{2.2cm}p{4.5cm}p{7cm}}
\toprule
\textbf{Metodo} & \textbf{Endpoint} & \textbf{Funcion} \\
\midrule
\endhead
\texttt{GET} & \texttt{/api/health} & Verifica estado del servicio. \\
\texttt{GET} & \texttt{/api/security-rules} & Devuelve reglas cargadas desde YAML. \\
\texttt{GET} & \texttt{/api/ai-attack-examples} & Devuelve ejemplos educativos de ataques IA. \\
\texttt{POST} & \texttt{/api/chat} & Procesa el mensaje y devuelve respuesta, traza, seguridad, tokens, intencion, agente y herramienta. \\
\bottomrule
\end{longtable}

\subsection{Pipeline implementado}

\begin{enumerate}
  \item \texttt{raw\_input}
  \item \texttt{normalization}
  \item \texttt{tokenization}
  \item \texttt{entity\_extraction}
  \item \texttt{intent\_detection}
  \item \texttt{intent\_reinforcement}
  \item \texttt{policy\_guard}
  \item \texttt{orchestrator}
  \item \texttt{agent\_selection}
  \item \texttt{tool\_execution}
  \item \texttt{response\_generation}
  \item \texttt{audit\_log}
\end{enumerate}

Cada etapa genera los campos \texttt{stage}, \texttt{status}, \texttt{input\_type}, \texttt{output\_type}, \texttt{input}, \texttt{output} y \texttt{metadata}.

\section{Modulo OpenIA/OpenAI-compatible}

Se agrego el modulo \texttt{openia\_intent\_reinforcer.py}. Su funcion es reforzar la intencion detectada usando una capa semantica local. No llama a APIs externas y no requiere \texttt{OPENAI\_API\_KEY}.

\subsection{Comportamiento}

\begin{itemize}
  \item Si la intencion base es \texttt{unknown\_intent} o \texttt{general\_question}, puede reemplazarla por una intencion mas especifica cuando encuentra senales suficientes.
  \item Si la intencion base ya es especifica, puede conservarla y elevar la confianza.
  \item Si no encuentra senales utiles, conserva la intencion original.
  \item Si se configura un modo externo no soportado, no realiza llamadas de red y reporta que el modo externo no esta habilitado.
\end{itemize}

\subsection{Ejemplo validado}

\begin{lstlisting}[language={},caption={Refuerzo de intencion}]
Entrada: tokens
Intencion base: unknown_intent
Intencion final: explain_tokens
Refuerzo aplicado: true
Etapas de traza: 12
\end{lstlisting}

\section{Frontend}

El frontend esta organizado en tres paneles:

\begin{itemize}
  \item Panel izquierdo: chat.
  \item Panel central: flujo React Flow y metricas.
  \item Panel derecho: inspector tecnico.
\end{itemize}

La interfaz muestra conversacion, nodos del pipeline, estado de cada etapa, JSON de entrada y salida, conteo estimado de tokens, costo ilustrativo, politica de seguridad, refuerzo OpenIA, agente seleccionado, herramienta ejecutada, reglas activadas y ejemplos defensivos de ataques IA.

\section{Seguridad implementada}

El sistema incorpora controles de seguridad basicos:

\begin{itemize}
  \item Validacion Pydantic para requests.
  \item Rechazo de mensajes vacios.
  \item Limite de longitud del mensaje.
  \item Deteccion simple de prompt injection.
  \item Deteccion de texto con forma de secreto.
  \item Redaccion de secretos en trazas.
  \item Allowlist de herramientas.
  \item Bloqueo fail closed ante politica insegura.
  \item Sin ejecucion arbitraria de codigo.
  \item Sin endpoints destructivos.
  \item Sin uso de claves reales.
  \item Sin servicios pagos ni APIs externas.
\end{itemize}

\section{Ejemplos de ataques IA incluidos}

El endpoint \texttt{/api/ai-attack-examples} devuelve ejemplos educativos para explicar riesgos y defensas.

\begin{longtable}{p{4cm}p{5cm}p{5cm}}
\toprule
\textbf{Categoria} & \textbf{Ejemplo resumido} & \textbf{Defensa esperada} \\
\midrule
\endhead
Prompt injection directa & Ignorar instrucciones anteriores. & Detectar bypass y bloquear. \\
Extraccion de prompt interno & Pedir prompts o politicas internas. & No revelar configuracion interna. \\
Abuso de herramientas & Forzar herramientas no autorizadas. & Aplicar allowlist. \\
Inyeccion indirecta & Documento externo intenta dar instrucciones. & Tratar contenido externo como datos. \\
Secreto en entrada & Usuario pega token o API key. & Redactar y bloquear segun politica. \\
Accion destructiva & Saltar confirmacion para eliminar datos. & Requerir confirmacion y permisos. \\
\bottomrule
\end{longtable}

\section{Evidencia de validacion}

Se ejecutaron validaciones de backend y frontend.

\subsection{Backend}

\begin{lstlisting}[language=bash,caption={Comandos de validacion backend}]
python -m compileall -q app
pip check
\end{lstlisting}

Resultado:

\begin{lstlisting}[language={},caption={Resultado de dependencias Python}]
No broken requirements found.
\end{lstlisting}

Pruebas funcionales:

\begin{lstlisting}[language={},caption={Pruebas funcionales backend}]
/api/health -> OK
/api/ai-attack-examples -> 6 ejemplos
/api/chat con "tokens" -> explain_tokens, refuerzo aplicado, 12 etapas
/api/chat con prompt injection -> bloqueado, 12 etapas
\end{lstlisting}

\subsection{Frontend}

\begin{lstlisting}[language=bash,caption={Comandos de validacion frontend}]
npm run build
npm audit --audit-level=moderate
\end{lstlisting}

Resultado:

\begin{lstlisting}[language={},caption={Resultado frontend}]
TypeScript compila sin errores.
Vite genera bundle de produccion.
0 vulnerabilities.
\end{lstlisting}

\section{Limitaciones del demo}

\begin{itemize}
  \item No usa modelos reales.
  \item No llama a OpenAI ni a otro proveedor externo.
  \item El conteo de tokens es aproximado.
  \item El costo es ilustrativo.
  \item El detector de intenciones es determinista.
  \item No persiste auditorias en base de datos.
  \item No implementa autenticacion ni roles reales.
\end{itemize}

\section{Posibles extensiones}

\begin{itemize}
  \item Integrar OpenAI real mediante una bandera explicita y \texttt{OPENAI\_API\_KEY} en variable de entorno.
  \item Validar salidas de modelo con esquemas estrictos.
  \item Agregar pruebas automatizadas con \texttt{pytest}.
  \item Persistir auditoria redacted.
  \item Agregar autenticacion y RBAC.
  \item Sustituir el estimador por un tokenizer real.
  \item Exportar trazas como JSON o PDF.
  \item Agregar vista de comparacion entre intencion base e intencion reforzada.
\end{itemize}

\section{Conclusiones}

El demo cumple el objetivo de mostrar de forma educativa lo que ocurre internamente cuando un chatbot procesa un mensaje. La arquitectura separa frontend, backend, reglas, agentes, herramientas, seguridad y trazabilidad.

La aplicacion queda lista para demostraciones locales y como base para futuras integraciones reales, manteniendo por defecto un modo seguro, offline y sin secretos.

\end{document}
